% titel: Gravitationswellen von zwei ubereinander kreisenden Massen
% skript: Kap. Gravitation, Virialsatz

Zwei gleiche Massen $m$ kreisen mit Winkelgeschwindigkeit $\omega$ im Abstand
$r$ um ihren gemeinsamen Schwerpunkt. Aus dem Kräftegleichgewicht erhält man
den Zusammenhang $r(\omega)$; mit Energie- und Virialsatz folgt die
mechanische Gesamtenergie $E_\mathrm{tot}$, und über die abgestrahlte
Leistung $P \propto \omega^{10/3}$ ergibt sich, warum Gravitationswellen
erst spät im Signal stark werden.

\begin{werkzeugkasten}[Werkzeuge]
\textbf{Kraftbilanz} auf einer Kreisbahn (Masse $m$ um Abstand $r/2$ vom
Schwerpunkt):
\begin{equation}
  G\frac{m^2}{r^2} = m\,\omega^2\,\frac r2
  \;\Rightarrow\;
  r^3 = \frac{2Gm}{\omega^2}\;.
\end{equation}

\textbf{Potentielle Energie} durch Integration der Gravitationskraft vom
Unendlichen:
\begin{equation}
  E_\mathrm{pot} = -G\frac{m^2}{r}\;,
\end{equation}

\textbf{Virialsatz} für ein gebundenes System ($\langle E_\mathrm{kin}\rangle
$ gegen $\langle E_\mathrm{pot}\rangle$ zum Mittelwert):
\begin{equation}
  E_\mathrm{kin} = -\frac12 E_\mathrm{pot} \;.
\end{equation}

\textbf{Abgestrahlte Leistung} (quadrupol-Näherung, gegeben):
\begin{equation}
  P = \frac{8G}{5c^5}\,m^2\, r^4 \omega^6\;.
\end{equation}

\textbf{Konstanten:} $c=\qty{3e8}{m/s}$, $G=\qty{6.674e-11}{m^3/(kg\,s^2)}$,
$M_\odot=\qty{1.989e30}{kg}$; GW-Frequenz ist die doppelte Umlauffrequenz.
\end{werkzeugkasten}

\begin{werkzeugkasten}[Vorgehen in allgemeiner Form]
(a) Kräftegleichgewicht. (b) Frequenz aus dem Signal ablesen, $r$ und Dichte
berechnen. (c) Energie durch Integration. (d) $E_\mathrm{tot}=E_\mathrm{kin}+
E_\mathrm{pot}$ und Virialsatz prüfen, $E_\mathrm{tot}(\omega)$. (e)
Vorzeichen von $\mathrm d\omega/\mathrm dt$. (f) $P(\omega)$.
\end{werkzeugkasten}

\begin{achtung}[Typische Fallen]
\begin{itemize}
  \item Jede Masse dreht im Radius $r/2$ (nicht $r$): Der Abstand $r$ ist die
       Distanz \emph{zwischen} den Massen.
  \item Gravitationswellen-Frequenz ist die doppelte Umlauffrequenz.
  \item $\omega$ variiert; es bei der Antwort auf das Abnehmen des Abstands
       beachten.
\end{itemize}
\end{achtung}
